\section{Tight Complexity of the LRALM Family}
\label{sec:algorithm-framework}

In this section, we resolve \hyperlink{Q1}{\textbf{(Q1)}} by establishing tight
iteration-complexity bounds for single-loop linearized ALMs on
Problem~\eqref{prob:main}. Taking alternating GDA steps on \(\mathcal P_s\)
in~\eqref{eq:potential} together with moving primal and dual centers gives the
LRALM update below.
 
\begin{definition}
\label{def:splalm}
Let \(\theta=(\rho,s,\gamma,\eta,\tau)\) satisfy \(\rho>0\),
\(0<\eta\le\gamma<\infty\), and \(0<\tau\le s<\infty\), and let
\((\alpha,\beta)\in[0,1]^2\). Starting from
\(\bu^0=(\bx^0,\blmbd^0,\bz^0,\bmu^0)\in\RR^n\times\RR^m\times\RR^n\times\RR^m\),
the method \(\mathrm{LRALM}_{\theta,\alpha,\beta}\) updates on \(\mathcal P_s\)
in~\eqref{eq:potential}, for \(k\ge0\),
\begin{equation}
\begin{aligned}
    \bx^{k+1}&=\bx^k-\eta\,\nabla_{\bx}\mathcal P_s(\bx^k,\blmbd^k;\bz^k,\bmu^k),\\
    \blmbd^{k+1}&=\blmbd^k+\tau\,\nabla_{\blmbd}\mathcal P_s(\bx^{k+1},\blmbd^k;\bz^k,\bmu^k),\\
    \bz^{k+1}&=\bz^k+\alpha(\bx^{k+1}-\bz^k),\\
    \bmu^{k+1}&=\bmu^k+\beta(\blmbd^{k+1}-\bmu^k), 
\end{aligned}
    \label{eq:splalm-updates}
\end{equation}
and outputs
\((\bx^{k+1},\by^{k+1},\bq^k)
:=(\bx^{k+1},\by(\bx^{k+1},\blmbd^k),\bq(\bx^{k+1},\blmbd^k))\), where
\(\by(\cdot,\cdot)\) and \(\bq(\cdot,\cdot)\) are defined
in~\eqref{eq:prox-multiplier}.
\end{definition} 

When applied to \(\calD_\rho\), existing single-loop methods from both ALM
and smooth minimax optimization take the form~\eqref{eq:splalm-updates},
with different choices of the dual regularizer \(s\) and the center
stepsizes \((\alpha,\beta)\). We first clarify the effect of \(s\) on
\(\calD_\rho\) and then organize these center-movement choices into the LRALM
family.

\subsection{The role of the dual regularizer and the LRALM family}
\label{subsec:lralm-family}

In linearized ALM~\eqref{eq:alm-reduced} on \(\calD_\rho\),
constraint violation is typically controlled by choosing a large penalty
\(\rho\)~\cite{luSingleLoopGradientDescent2022,sun2024dual_descent_alm_admm,elbourkhissi2025perturbed}.
Indeed, combining~\eqref{eq:reduced-al-gradient} with the bound
\(\|\bq(\bx,\blmbd)\|\le L_h\) from Lemma~\ref{lem:reduced-al}(a) and
Assumption~\ref{ass:h} gives
\[
    \|\calA(\bx)-\by(\bx,\blmbd)\|
    \le\frac{L_h+\|\blmbd\|}{\rho}.
\]
On the other hand, since \(L_{\calD,\bx}\) grows linearly with \(\rho\),
a large \(\rho\) increases the primal smoothness modulus in
Proposition~\ref{prop:curvature-smoothness}(a) and restricts the primal stepsize
\(\eta\). The next lemma shows that,
at a fixed \(\rho\), maximizing over \(\blmbd\) with the dual regularizer
raises the effective penalty from \(\rho\) to \(\rho+s\). The identity
follows from~\eqref{eq:dual-smoothing}, since the infimal convolution of
\(\|\cdot\|^2/(2\rho)\) and \(\|\cdot\|^2/(2s)\) is
\(\|\cdot\|^2/(2(\rho+s))\); the full proof is given in
\supploc{app:effective-penalty}.

\begin{lemma}
\label{lem:effective-penalty}
Suppose that Assumption~\ref{ass:h} holds, and let \(\rho,s>0\). For every
\((\bx,\bmu)\in\RR^n\times\RR^m\),
\begin{equation}
    \max_{\blmbd\in\RR^m}\Bigl\{\calD_\rho(\bx,\blmbd)
    -\frac{1}{2s}\|\blmbd-\bmu\|^2\Bigr\}
    =\calD_{\rho+s}(\bx,\bmu).
    \label{eq:effective-penalty}
\end{equation}
Moreover, the left-hand side of~\eqref{eq:effective-penalty} has a unique
maximizer \(\blmbd^\star(\bx,\bmu)\), which satisfies
\[
    \by\bigl(\bx,\blmbd^\star(\bx,\bmu)\bigr)
    =\prox_{h/(\rho+s)}\bigl(\calA(\bx)+\bmu/(\rho+s)\bigr),
    \qquad
    \bq^\star:=\bq\bigl(\bx,\blmbd^\star(\bx,\bmu)\bigr)
    \in\partial h\bigl(\by(\bx,\blmbd^\star(\bx,\bmu))\bigr),
\]
and
\begin{equation}
\begin{aligned}
    \blmbd^\star(\bx,\bmu)
    &=\frac{s}{\rho+s}\bq^\star+\frac{\rho}{\rho+s}\bmu,\\
    \nabla_{\bmu}\calD_{\rho+s}(\bx,\bmu)
    &=\calA(\bx)-\by\bigl(\bx,\blmbd^\star(\bx,\bmu)\bigr)
    =\frac{\bq^\star-\bmu}{\rho+s}
    =\frac{\blmbd^\star(\bx,\bmu)-\bmu}{s}.
\end{aligned}
    \label{eq:effective-penalty-maximizer}
\end{equation}
\end{lemma}

Under Assumption~\ref{ass:h}, we have \(\|\bq^\star\|\le L_h\),
so~\eqref{eq:effective-penalty-maximizer} replaces \(\rho\) in the previous
bound by \(\rho+s\) and gives
\[
    \bigl\|\calA(\bx)-\by\bigl(\bx,\blmbd^\star(\bx,\bmu)\bigr)\bigr\|
    \le\frac{L_h+\|\bmu\|}{\rho+s}.
\]
Meanwhile, the primal smoothness modulus \(L_x=\gamma^{-1}+L_{\calD,\bx}\)
of \(\mathcal P_s\) is independent of \(s\) by
Proposition~\ref{prop:curvature-smoothness}(a). Increasing \(s\) therefore
strengthens the effective penalty without shrinking the primal stepsize
\(\eta\).

Whether \(s\) enters the dual update of~\eqref{eq:splalm-updates}, however, depends
on the movement of the dual center \(\bmu^k\). In the dual step
of~\eqref{eq:splalm-updates}, the dual regularizer contributes
\(-(\blmbd^k-\bmu^k)/s\) to
\(\nabla_{\blmbd}\mathcal P_s(\bx^{k+1},\blmbd^k;\bz^k,\bmu^k)\). When
\(\beta=1\), we have \(\bmu^k=\blmbd^k\) for all \(k\ge1\), so this
regularization term vanishes thereafter and the update reverts to the penalty
\(\rho\). When \(\beta<1\), the dual regularizer remains active; in particular,
taking \(\beta=0\) with \(\bmu^0=\zero\) fixes \(\bmu^k\equiv\zero\), matching
the effective penalty \(\rho+s\) of Lemma~\ref{lem:effective-penalty}. The
parameter \(\alpha\in[0,1]\) plays the analogous role for the primal center
\(\bz^k\), from a fixed center at \(\alpha=0\) to \(\bz^k=\bx^k\) for all
\(k\ge1\) at \(\alpha=1\). For a given
parameter tuple \(\theta\), collecting all choices of
\((\alpha,\beta)\in[0,1]^2\) defines the LRALM family, whose worst-case
iteration complexity over
\(\mathcal I(\mathbf C,\Delta)\)
for Problem~\eqref{prob:main} is defined below.

\begin{figure}[t]
\centering
\definecolor{cSlate}{RGB}{65,78,96}
\definecolor{cRed}{RGB}{178,56,50}
\definecolor{cTeal}{RGB}{20,118,92}
\definecolor{cNavy}{RGB}{32,74,158}
\resizebox{\ifdim\width>\textwidth\textwidth\else\width\fi}{!}{%
\begin{tikzpicture}[x=5.6cm,y=4.4cm,font=\small,
  axis/.style={cSlate,line width=0.85pt,-{Stealth[length=4.5pt]}},
  mod/.style={cSlate!85,line width=0.9pt,-{Stealth[length=4.8pt]}}]
  % region (0,1)^2
  \fill[cNavy!8] (0,0) rectangle (1,1);
  \draw[cSlate!40,dashed,line width=0.75pt] (1,0)--(1,1);
  % coordinate axes
  \draw[axis] (0,-0.02)--(0,1.20) node[above,text=cSlate] {$\beta$\ \footnotesize(dual center $\bmu^k$)};
  \draw[axis] (-0.02,0)--(1.15,0) node[below=3pt,text=cSlate,align=center] {$\alpha$\\[-1pt]\footnotesize(primal center $\bz^k$)};
  % highlighted 1D edges: beta=1 (red) and beta=0 (teal)
  \draw[cRed,line width=2.3pt] (0,1)--(1,1);
  \draw[cTeal,line width=2.3pt] (0,0)--(1,0);
  % open endpoints at alpha=0, solid corner points at alpha=1
  \node[circle,draw=cRed,fill=white,line width=1.1pt,inner sep=2.0pt] at (0,1) {};
  \node[circle,draw=cTeal,fill=white,line width=1.1pt,inner sep=2.0pt] at (0,0) {};
  \node[circle,draw=white,fill=cRed,line width=0.8pt,inner sep=2.7pt] at (1,1) {};
  \node[circle,draw=white,fill=cTeal,line width=0.8pt,inner sep=2.7pt] at (1,0) {};
  % axis labels
  \node[below,text=cSlate] at (0,-0.04) {$0$};
  \node[below,text=cSlate] at (1,-0.04) {$1$};
  \node[left,align=right] at (-0.04,1) {\textcolor{cSlate}{$1$}\\[-1pt]
    \footnotesize\textcolor{cRed}{$\bmu^k=\blmbd^k$: penalty $\rho$, $s$ inactive}};   
  \node[left,align=right] at (-0.04,0) {\textcolor{cSlate}{$0$}\\[-1pt]
    \footnotesize\textcolor{cTeal}{$\bmu^k\equiv\zero$: penalty $\rho+s$}};
  % methods colored by geometric region
  \node[anchor=north,text=cRed] at (0.5,0.96)
    {Smoothed GDA};
  \node[align=center,text=cNavy] at (0.5,0.5)
    {DS-GDA\\[-1pt] 
    \scriptsize\textcolor{cNavy!75}{$(\alpha,\beta)\in(0,1)^2$}};
  \node[anchor=south,text=cTeal] at (0.5,0.04)
    {Perturbed Smoothed GDA};
  \node[cRed,anchor=south west,align=left] at (1.03,0.99)
    {Linearized ALM\\[-1pt]
    \scriptsize\textcolor{cRed!80}{(GDA on $\calD_\rho$, $\tau=\rho$)}};
  \node[cTeal,anchor=south west] at (1.03,0.04)
    {Perturbed GDA};
  % the two regularization directions
  \draw[mod] (0.91,1.09)--(0.09,1.09)
     node[midway,above=1.5pt,font=\footnotesize,text=cSlate]
     {primal smooth.\ $+\frac{1}{2\gamma}\|\bx-\bz^k\|^2$};
  \draw[mod] (1.08,0.90)--(1.08,0.20)
     node[midway,right=2pt,align=left,font=\footnotesize,text=cSlate]
     {dual reg.\\[-1pt]$-\frac{1}{2s}\|\blmbd-\bmu^k\|^2$};
\end{tikzpicture}% 
}
\caption{Center-stepsize space \((\alpha,\beta)\in[0,1]^2\) of the LRALM family 
\(\calC_\theta\) acting on \(\calD_\rho\). Red, teal, and blue mark the top
edge \(\beta=1\), the bottom edge \(\beta=0\) (with \(\bmu^0=\zero\)), and the
region \((0,1)^2\), while solid circles mark the corner methods at
\(\alpha=1\) and open circles mark the excluded endpoints at \(\alpha=0\).} 
\label{fig:lralm-family}
\end{figure} 

\begin{definition}
\label{def:complexity-measure}
Given \(\theta=(\rho,s,\gamma,\eta,\tau)\) as in Definition~\ref{def:splalm},
the LRALM family is
\(\calC_\theta:=\{\mathrm{LRALM}_{\theta,\alpha,\beta}:(\alpha,\beta)\in[0,1]^2\}\).
Given \(\mathbf C\in\RR^5_{\ge0}\), define the admissible initialization set
\[  
    \mathcal S_0(\mathbf C)
    :=\bigl\{\bu^0=(\bx^0,\blmbd^0,\bz^0,\bmu^0):\
        \bx^0=\zero,\ \bmu^0=\zero,\ \|\blmbd^0\|\le L_h,\ \|\bz^0\|\le1\bigr\}.
\]
For \(\epsilon>0\), \(I\in\mathcal I(\mathbf C,\Delta)\), and
\(\bu^0\in\mathcal S_0(\mathbf C)\), the iteration complexity of 
\(\mathrm{LRALM}_{\theta,\alpha,\beta}\) on \((I,\bu^0)\) is
\[
    T_\epsilon(\theta,\alpha,\beta;I,\bu^0) 
    :=\inf\Bigl\{K\ge1:\ \min_{0\le k<K}
        \mathcal R_{\mathrm{KKT}}(\bx^{k+1},\by^{k+1},\bq^k)\le\epsilon\Bigr\},
\]
with \(\inf\varnothing=+\infty\). The worst-case iteration complexity of
\(\mathrm{LRALM}_{\theta,\alpha,\beta}\) over \(\mathcal I(\mathbf C,\Delta)\) is
\[
    T_\epsilon\bigl(\theta,\alpha,\beta;\mathcal I(\mathbf C,\Delta)\bigr)
    :=\sup_{\substack{I\in\mathcal I(\mathbf C,\Delta)\\ \bu^0\in\mathcal S_0(\mathbf C)}}\
    T_\epsilon(\theta,\alpha,\beta;I,\bu^0),
\] 
and that of \(\calC_\theta\) is
\[
    T_\epsilon\bigl(\calC_\theta,\mathcal I(\mathbf C,\Delta)\bigr)
    :=\inf_{(\alpha,\beta)\in[0,1]^2}\
    T_\epsilon\bigl(\theta,\alpha,\beta;\mathcal I(\mathbf C,\Delta)\bigr).
\]
\end{definition} 

Figure~\ref{fig:lralm-family} locates existing single-loop methods applied
to \(\calD_\rho\) in the \((\alpha,\beta)\) plane of \(\calC_\theta\). On the
top edge \(\beta=1\), \(\alpha=1\) with \(\tau=\rho\) recovers linearized
ALM~\eqref{eq:alm-reduced} and \(\alpha\in(0,1)\) gives
Smoothed GDA~\cite{zhang2020smoothed}; on the bottom edge \(\beta=0\)
with \(\bmu^0=\zero\), \(\alpha=1\) yields Perturbed
GDA~\cite{lu2020hybrid,xu2023unified} and \(\alpha\in(0,1)\) yields
Perturbed Smoothed GDA~\cite{li2026smoothing}. The interior
\((\alpha,\beta)\in(0,1)^2\) corresponds to doubly smoothed GDA
(DS-GDA)~\cite{zheng2023universal}.

\subsection{\texorpdfstring{Tight bounds in two effective-penalty regimes}{Tight bounds in two effective-penalty regimes}}
\label{sec:lower-bound}

We now establish matching lower and upper bounds for \(\calC_\theta\) in
two regimes of the effective penalty \(\rho+s\); the proofs are given in
\ifsupp Appendices~\ref{app:lower-bound} and~\ref{app:upper-bounds}\else
the supplementary material\fi. By reusing
\(\mathcal O_{f,\calA}(\bx^{k+1})\) in the next primal step, each iteration
of~\eqref{eq:splalm-updates} makes one call to \(\mathcal O_{f,\calA}\) and
two calls to \(\mathcal O_h\), so the iteration bounds below agree with the
oracle complexity up to a constant factor.

\subsubsection{Tight bounds for \texorpdfstring{\(\rho+s\le c_1/\epsilon\)}{rho+s <= c1/epsilon}} 

We first establish that, while \(\epsilon(\rho+s)\) stays bounded by a
constant, \(\calC_\theta\) attains \(\epsilon^{-4}\) and no
center-stepsize pair \((\alpha,\beta)\in[0,1]^2\) improves on it.

\begin{theorem}
\label{thm:lower-bound}
Let \(\mathbf C\in\RR^5_{\ge0}\) satisfy \(L_f,G_f,G_{\calA},L_h\ge1\),
and let \(0<\eta\le\gamma\), \(\tau>0\), and \(\Delta>0\). Then there exist
constants \(c_1,c_2>0\), independent of \((\rho,s)\), such
that, for every \(\epsilon\in(0,c_2)\), every \(\rho>0\), and every
\(s\ge\tau\) with \(\epsilon(\rho+s)\le c_1\),
\[
    T_\epsilon\bigl(\calC_\theta,\mathcal I(\mathbf C,\Delta)\bigr)
    =\Omega(\epsilon^{-4}),
\]
where the constant in \(\Omega(\cdot)\) is independent of \((\rho,s)\).
\end{theorem}
The matching upper bound is attained by Smoothed GDA~\cite{zhang2020smoothed}, 
the member of \(\calC_\theta\) with \(\beta=1\). Since
\(\bmu^k=\blmbd^k\) for all \(k\ge1\) at \(\beta=1\), the LRALM updates
from \(k=1\) on coincide with Smoothed GDA and no longer involve \(s\), and the
bound follows
from~\cite[Theorem~B.1]{li2026smoothing} with \(\alpha=c\,\epsilon^2\); the
full proof is given in \supploc{app:upper-beta1}.

\begin{proposition}
\label{prop:upper-beta1}
Let \(\mathbf C\in\RR^5_{\ge0}\), \(\Delta>0\), and \(\rho>0\), set
\(\ell:=\max\{L_{\calD,\bx},G_{\calA},1/\rho\}\), and let
\(0<\gamma\le\bar c/\ell\), \(0<\eta\le\gamma/2\), and
\(0<\tau\le\bar c(\eta/\gamma)^2/\ell\) for a sufficiently small
constant \(\bar c>0\). Then there exists \(c>0\) such that, for
every \(s\ge\tau\) and every \(\epsilon\in(0,1]\), the choice
\((\alpha,\beta)=(c\,\epsilon^2,1)\) yields
\[
    T_\epsilon\bigl(\theta,c\,\epsilon^2,1;\mathcal I(\mathbf C,\Delta)\bigr)
    =\mathcal O(\epsilon^{-4}),
\]
where the constant in \(\mathcal O(\cdot)\) is independent of \(s\).
\end{proposition}
\noindent Under the stepsize conditions of
Proposition~\ref{prop:upper-beta1}, combining Theorem~\ref{thm:lower-bound}
and Proposition~\ref{prop:upper-beta1}
yields \(T_\epsilon(\calC_\theta,\mathcal I(\mathbf C,\Delta))=\Theta(\epsilon^{-4})\)
whenever \(\rho+s\le c_1/\epsilon\). Since \(c_1\) and \(c_2\) are
independent of \((\rho,s)\), improving on \(\Theta(\epsilon^{-4})\) over
\(\mathcal I(\mathbf C,\Delta)\) requires \(\rho+s>c_1/\epsilon\) for 
all sufficiently small \(\epsilon\). This partially 
answers~\hyperlink{Q1}{\textbf{(Q1)}}. Specifically, for a fixed-penalty linearized ALM with
\(\rho=\Theta(1)\), improving on \(\epsilon^{-4}\) requires \(s\) to grow as the
accuracy improves, so \(s\) cannot remain \(\Theta(1)\).

\subsubsection{Tight bounds for \texorpdfstring{\(s=c_s/\epsilon\)}{s = cs/epsilon}}

To reach the regime \(\rho+s>c_1/\epsilon\) with \(\rho>0\) fixed, we
scale \(s=c_s/\epsilon\). At this scale, \(\calC_{\theta(\epsilon)}\)
attains \(\epsilon^{-3}\), and no center-stepsize pair
\((\alpha,\beta)\in[0,1]^2\) improves on it.

\begin{theorem}
\label{thm:lower-bound-fixed-rho}
Let \(\mathbf C\in\RR^5_{\ge0}\) satisfy \(L_f,G_f,G_{\calA},L_h\ge1\),
and let \(\rho>0\), \(0<\eta\le\gamma\), \(\tau>0\), \(c_s>0\), and
\(\Delta>0\). Then there exist constants \(B\ge1\) and
\(\epsilon_0\in(0,1]\) such that, for every
\(\epsilon<\epsilon_0\) and every \((\alpha,\beta)\in[0,1]^2\), with
\(\theta(\epsilon):=(\rho,c_s/\epsilon,\gamma,\eta,\tau)\),
\[
    T_\epsilon\bigl(\theta(\epsilon),\alpha,\beta;\mathcal I(\mathbf C,\Delta)\bigr)
    \begin{cases}
        =+\infty, & \alpha=0\ \text{or}\ \alpha>B\epsilon,\\[2pt]
        =\Omega\bigl(\alpha^{-1}\epsilon^{-2}\bigr), & 0<\alpha\le B\epsilon.
    \end{cases}
\]
In particular,
\(T_\epsilon(\calC_{\theta(\epsilon)},\mathcal I(\mathbf C,\Delta))=\Omega(\epsilon^{-3})\).
\end{theorem}
\noindent The matching upper bound is attained by Perturbed Smoothed
GDA, the member of \(\calC_{\theta(\epsilon)}\) with
\(\beta=0\). At \(\beta=0\), \(\bmu^k\equiv\zero\) and LRALM coincides
with Perturbed Smoothed GDA with dual perturbation \(r_y=1/s\). A sufficiently
large \(c_s\) keeps the induced stationarity bias, of order \(1/s\), at most
\(\epsilon/4\), and~\cite[Condition~3.2(ii)]{li2026smoothing} admits
\(\alpha=c_\alpha\epsilon\) for a sufficiently small \(c_\alpha>0\), so
that~\cite[Theorem~3.3(i)]{li2026smoothing} yields the bound; the full proof
is given in \supploc{app:upper-beta0}.

\begin{proposition}
\label{prop:upper-beta0}
Let \(\mathbf C\in\RR^5_{\ge0}\), \(\Delta>0\), and \(\rho>0\), let
\(\ell,\gamma,\eta,\tau\) be as in Proposition~\ref{prop:upper-beta1}, and let
\(c_s>0\) be sufficiently large. Then there exist constants \(c_\alpha>0\)
and \(\bar\epsilon\in(0,1]\) such that, for every \(\epsilon\le\bar\epsilon\), with
\(\theta(\epsilon):=(\rho,c_s/\epsilon,\gamma,\eta,\tau)\), the choice
\((\alpha,\beta)=(c_\alpha\epsilon,0)\) yields
\[
    T_\epsilon\bigl(\theta(\epsilon),c_\alpha\epsilon,0;\mathcal I(\mathbf C,\Delta)\bigr)
    =\mathcal O(\epsilon^{-3}).
\]
\end{proposition}
\noindent Combining Theorem~\ref{thm:lower-bound-fixed-rho} and
Proposition~\ref{prop:upper-beta0} yields the tight family complexity
\(T_\epsilon(\calC_{\theta(\epsilon)},\mathcal I(\mathbf C,\Delta))
=\Theta(\epsilon^{-3})\) for every sufficiently large \(c_s\).

Taken together, the two regimes answer~\hyperlink{Q1}{\textbf{(Q1)}}.
For a fixed-penalty linearized ALM with \(\rho=\Theta(1)\), \(s\) must grow at
least as \(\epsilon^{-1}\) to improve on \(\epsilon^{-4}\), and
\(s=c_s/\epsilon\) with a sufficiently large \(c_s\) suffices for
\(\epsilon^{-3}\). Within \(\calC_{\theta(\epsilon)}\), however,
\(\epsilon^{-3}\) cannot be improved, and the limitation comes from the primal
center stepsize \(\alpha\), that is, from how the primal center moves,
independently of the dual center. This motivates 
the accelerated method in Section~\ref{sec:acceleration}.

\begin{remark} 
\begin{enumerate}[label=(\roman*)]
\item Recent lower bounds establish the optimal first-order complexity for
general smooth NC-C minimax problems~\cite{pan2026optimal,zhang2026matching}
and certify the tightness of individual single-loop methods under prescribed
stepsize
conditions~\cite[Theorem~3.4(i) and Corollary~3.1(i)]{li2026smoothing}.
Complementing these results, Theorems~\ref{thm:lower-bound}
and~\ref{thm:lower-bound-fixed-rho} establish lower bounds on the structured
minimax problems \(\min_{\bx}\max_{\blmbd}\calD_\rho(\bx,\blmbd)\)
generated by \(\mathcal I(\mathbf C,\Delta)\) across every constant
center-stepsize pair \((\alpha,\beta)\in[0,1]^2\). Specifically, because
these problems are smooth NC-C minimax problems, these bounds transfer
directly to the corresponding single-loop methods on
general smooth NC-C problems.
\item In nonconvex constrained optimization, existing single-loop ALM analyses
typically establish complexity guarantees through growing
penalties~\cite{luSingleLoopGradientDescent2022,sun2024dual_descent_alm_admm,elbourkhissi2025perturbed}.
Alongside these developments, combining Theorems~\ref{thm:lower-bound}
and~\ref{thm:lower-bound-fixed-rho} with Propositions~\ref{prop:upper-beta1}
and~\ref{prop:upper-beta0} establishes tight worst-case complexity bounds for
the fixed-penalty LRALM family \(\calC_\theta\). Specifically,
the effective-penalty threshold shows that the required \(\epsilon^{-1}\)
scaling can be carried by \(s\) while \(\rho\) stays fixed, so the primal
smoothness modulus and stepsize remain constant.
\end{enumerate}
\end{remark}

\subsection{Proof sketch of the lower bounds}
\label{sec:lb-proof}
We sketch the proofs of Theorems~\ref{thm:lower-bound}
and~\ref{thm:lower-bound-fixed-rho} in two steps; the explicit constructions of
Instances~II and~IV and all proofs are given in \supploc{app:lower-bound}.
Throughout this subsection, we write
\(R_k:=\mathcal R_{\mathrm{KKT}}(\bx^{k+1},\by^{k+1},\bq^k)\) and assume
\(L_f,G_f,G_{\calA},L_h\ge1\).

\textit{Step~1.} We partition \([0,1]^2\) by the primal center stepsize
\(\alpha\), with threshold \(\alpha_\epsilon=\Theta(\epsilon^2)\) for
\(\rho+s\le c_1/\epsilon\) and \(B\epsilon\) for \(s=c_s/\epsilon\). Below
either threshold, the plateau Instance~I forces
\(\Omega(\Delta/(\gamma\alpha\epsilon^2))\) iterations for every \(\beta\).

\textit{Step~2.} Above these thresholds, Instances~II, III, and~IV yield
\(T_\epsilon=+\infty\). Under \(\rho+s\le c_1/\epsilon\), the rotating
Instance~II covers \(\beta\in(0,1]\) and the Huber Instance~III sustains a dual
feasibility bias at \(\beta=0\); under \(s=c_s/\epsilon\), the radial
Instance~IV covers every \(\beta\in[0,1]\). Taking the infimum over
\((\alpha,\beta)\in[0,1]^2\) of the bounds in the two steps yields both
theorems.

\subsubsection{Hard instances for \texorpdfstring{\(\rho+s\le c_1/\epsilon\)}{rho+s <= c1/epsilon}}

\paragraph{I. Plateau instance.}
The first instance covers \(\alpha\le\alpha_\epsilon\) for every
\(\beta\in[0,1]\), and is reused below for \(\alpha\le B\epsilon\).

\begin{example}
\label{ex:plateau}
Let \(0<\epsilon\le\min\{1/2,\sqrt{\Delta}/2\}\) and
\(x_\star:=\Delta/(2\epsilon)+\epsilon\). Define \(\calA\equiv0\),
\(h(y)=|y|\), and \(f\colon\RR\to\RR\) by
\[
f(x)=
\begin{cases}
2\epsilon(x_\star-x)-2\epsilon^2, & x<x_\star-2\epsilon,\\[4pt]
\dfrac12(x-x_\star)^2, & |x-x_\star|\le2\epsilon,\\[4pt]
2\epsilon(x-x_\star)-2\epsilon^2, & x>x_\star+2\epsilon.
\end{cases}
\]
\end{example}

Because \(\calA\equiv0\), the dual variables remain zero from \(\bu^0=\zero\),
and every iterate on \(x<x_\star-2\epsilon\) satisfies
\(R_k=|\nabla f(x^{k+1})|=2\epsilon\) and \(|x^k-z^k|\le2\gamma\epsilon\). A
small \(\alpha\) therefore limits each center step to order
\(\gamma\alpha\epsilon\), yielding the lower bound below; its proof is given in
\supploc{subsec:lb-plateau}.

\begin{lemma}
\label{lem:main-plateau}
Let \((\alpha,\beta)\in[0,1]^2\), and let \(\epsilon>0\) satisfy the
condition of Example~\ref{ex:plateau} and \(\epsilon^2\le\Delta/(16\gamma)\).
Then the instance \(I\) of Example~\ref{ex:plateau} belongs to
\(\mathcal I(\mathbf C,\Delta)\), and \(\mathrm{LRALM}_{\theta,\alpha,\beta}\)
initialized at \(\bu^0=\zero\) satisfies
\(T_\epsilon(\theta,\alpha,\beta;I,\zero)>\Delta/(16\gamma\alpha\epsilon^2)\)
for \(\alpha>0\), and \(T_\epsilon(\theta,\alpha,\beta;I,\zero)=+\infty\) for
\(\alpha=0\).
\end{lemma}

\paragraph{II. Rotating instance.}
For \(\alpha>\alpha_\epsilon\) and \(\beta\in(0,1]\), we build on the
quadratic coupling of~\cite[Example~5.1]{li2026smoothing} with a radial
extension that keeps \(f\) globally \(1\)-smooth and \(1\)-Lipschitz while
inducing a non-decaying unit-modulus mode in the linear LRALM recursion; the
explicit instance and proof are given in \ifsupp Example~\ref{ex:rotation} and
Appendix~\ref{subsec:lb-rotation}\else the supplementary material\fi.

\begin{lemma}
\label{lem:main-rotation}
Under the assumptions of Theorem~\ref{thm:lower-bound}, let
\(\alpha>\alpha_\epsilon\) and \(\beta\in(0,1]\). Then there exists an
instance \(I\in\mathcal I(\mathbf C,\Delta)\) on which
\(\mathrm{LRALM}_{\theta,\alpha,\beta}\) initialized at \(\bu^0=\zero\)
satisfies \(R_k\ge\frac76\epsilon\) for all \(k\ge0\); hence
\(T_\epsilon(\theta,\alpha,\beta;I,\zero)=+\infty\).
\end{lemma}

\paragraph{III. Huber instance.}
For \(\alpha>\alpha_\epsilon\) and \(\beta=0\), we retain \(f\) from
Example~\ref{ex:plateau} and pair a constant mapping \(\calA\) with a Huber
penalty \(h\).

\begin{example}
\label{ex:huber}
Let \(\epsilon\) and \(f\) be as in Example~\ref{ex:plateau}, and set
\[
q_h:=2s+1,
\qquad
a_0:=\frac{4\epsilon(\rho+3s+1)}{q_h}.
\]
Define \(\calA\equiv a_0\) and \(h\colon\RR\to\RR\) by
\[
h(y)=
\begin{cases}
\dfrac{q_h}{2}y^2, & |y|\le q_h^{-1},\\[4pt]
|y|-\dfrac{1}{2q_h}, & |y|>q_h^{-1}.
\end{cases}
\]
\end{example}

At \(\beta=0\), the dual center remains at \(\mu^k=0\), so the regularizer
\(-\frac{1}{2s}\lambda^2\) pulls \(\lambda^k\) toward zero and keeps it below
the multiplier \(q_ha_0\) needed for feasibility \(y^{k+1}=a_0\), so that
\(\lambda^k\) converges instead to \(sq_ha_0/(\rho+q_h+s)\). This bias leads to the following lower
bound, whose proof is deferred to \supploc{subsec:lb-huber}.

\begin{lemma}
\label{lem:main-huber}
Let \(\alpha\in[0,1]\) and \(\beta=0\), and suppose \(\epsilon\) satisfies
the condition of Example~\ref{ex:plateau} and \(\epsilon(\rho+s)\le c_1\).
Then the instance \(I\) of Example~\ref{ex:huber} belongs to
\(\mathcal I(\mathbf C,\Delta)\), and \(\mathrm{LRALM}_{\theta,\alpha,0}\)
initialized at \(\bu^0=\zero\) satisfies \(R_k\ge3\epsilon\) for all
\(k\ge0\); hence \(T_\epsilon(\theta,\alpha,0;I,\zero)=+\infty\).
\end{lemma}

Combining Lemmas~\ref{lem:main-plateau}--\ref{lem:main-huber} over the three
regions of \([0,1]^2\) yields Theorem~\ref{thm:lower-bound}, whose complete
proof appears in \supploc{subsec:lb-combine}.

\subsubsection{Hard instance for \texorpdfstring{\(s=c_s/\epsilon\)}{s = cs/epsilon}}

For \(\alpha\le B\epsilon\) and any \(\beta\in[0,1]\),
Lemma~\ref{lem:main-plateau} on Example~\ref{ex:plateau} directly gives the
lower bound \(\Omega(\Delta\gamma^{-1}\epsilon^{-3})\).

\paragraph{IV. Radial instance.}
For \(\alpha>B\epsilon\), rescaling the rotating construction of Instance~II and
initializing on its rotating eigenmode removes the transient term for every
\(\beta\in[0,1]\), keeping the iterates on a circle with non-decaying KKT
residuals; the explicit instance and proof are given in
\ifsupp Example~\ref{ex:radial} and Appendix~\ref{app:radial}\else the
supplementary material\fi.

\begin{lemma}
\label{lem:main-radial}
There exist constants \(B\ge1\), \(\epsilon_{\mathrm{r}}\in(0,\tfrac12]\), and
\(c_{\mathrm{r}}>0\) such that, for every \(\epsilon\in(0,\epsilon_{\mathrm{r}})\),
\(\alpha\in(B\epsilon,1]\), and \(\beta\in[0,1]\), there exist an
instance \(I\in\mathcal I(\mathbf C,\Delta)\) and an initialization
\(\bu^0\in\mathcal S_0(\mathbf C)\) such that
\(\mathrm{LRALM}_{\theta(\epsilon),\alpha,\beta}\) initialized at \(\bu^0\)
satisfies \(R_k\ge c_{\mathrm{r}}\alpha>\epsilon\) for all \(k\ge0\); hence
\(T_\epsilon(\theta(\epsilon),\alpha,\beta;I,\bu^0)=+\infty\).
\end{lemma}

Lemmas~\ref{lem:main-plateau} and~\ref{lem:main-radial} together yield
Theorem~\ref{thm:lower-bound-fixed-rho}, whose complete proof appears in
\supploc{subsec:lb-fixed-rho-proof}.
